\documentclass[12pt]{article}
\usepackage[utf8]{inputenc}
\usepackage{graphicx}
\usepackage{amsmath}
\usepackage{amssymb}
\usepackage{mathrsfs}
\usepackage{pgfornament}
\usepackage[many]{tcolorbox}
\usepackage[margin = 0.5in]{geometry}
\usepackage{collectbox}
\usepackage{mdframed}
\usepackage{xcolor}
\usepackage{mathtools}
\DeclarePairedDelimiter\ceil{\lceil}{\rceil}
\DeclarePairedDelimiter\floor{\lfloor}{\rfloor}

\newcommand\textbox[1]{%
  \parbox{.333\textwidth}{#1}%
}

\linespread{2}

\makeatletter
\newcommand{\mybox}{%
    \collectbox{%
        \setlength{\fboxsep}{2pt}%
        \fbox{\BOXCONTENT}%
    }%
}
\makeatother

\usepackage{listings}
\usepackage{color}

\definecolor{dkgreen}{rgb}{0,0.6,0}
\definecolor{gray}{rgb}{0.5,0.5,0.5}
\definecolor{mauve}{rgb}{0.58,0,0.82}

\lstset{frame=tb,
  language=R,
  aboveskip=3mm,
  belowskip=3mm,
  showstringspaces=false,
  columns=flexible,
  basicstyle={\small\ttfamily},
  numbers=none,
  numberstyle=\tiny\color{gray},
  keywordstyle=\color{blue},
  commentstyle=\color{dkgreen},
  stringstyle=\color{mauve},
  breaklines=true,
  breakatwhitespace=true,
  tabsize=3
}
\title{MAT 523 HW 5}
\begin{document}

\begin{center}
\begin{large}
\textbf{MAT 523: Functions of a Real Variable I \\
Homework V \\}
\end{large}
Nolan R. H. Gagnon \\
\end{center}

\noindent \textbf{(1.)}

\vspace{3mm}

\noindent\mybox{\colorbox{lightgray}{\textbf{Statement}}}\hspace{3mm}\hrulefill

\noindent Every interval (open, closed, half-open, bounded, unbounded, etc.) is a Borel set.

\vspace{5mm}

\noindent\mybox{\colorbox{lightgray}{\textbf{Proof}}}\hspace{3mm}\hrulefill

The collection $\mathcal{B}$ of Borel sets is the smallest $\sigma-$algebra of sets of real numbers that contain all the open sets of real numbers. Since open intervals are open sets, they are Borel sets. Let $a$ be a finite real number. From our previous conclusion, we know that $(a,\infty)\in\mathcal{B}$ (since it is an open interval). But $\mathcal{B}$ is a $\sigma-$algebra, so it is closed with respect to the formation of complements. Thus, $(-\infty, a]\in\mathcal{B}.$ By similar reasoning, we can conclude that intervals of the form $[a, \infty)$ are members of $\mathcal{B}$. Next, since $\mathcal{B}$ is closed under countable unions, sets of the form $(-\infty,a)\cup(a,\infty)$ are members of $\mathcal{B}.$ But then complements of these types of sets are also in $\mathcal{B}$, so
\begin{align*}
    \left((-\infty,a)\cup(a,\infty)\right)^{C} &= [a,\infty)\cap(-\infty, a] \\
    &= \lbrace a \rbrace
\end{align*}
is in $\mathcal{B}.$ Using this result, we can show that closed intervals are members of $\mathcal{B}.$ Let $b$ be another finite real number. Then
\begin{align*}
    [a,b] &= (a,b)\cup\lbrace a \rbrace \cup \lbrace b \rbrace 
\end{align*}
is in $\mathcal{B}$ because it is a countable union of sets that we already know to be in $\mathcal{B}.$ By similar reasoning, we can show that bounded, half-open intervals are in $\mathcal{B}.$

\hfill$\blacksquare$

\pagebreak

\noindent \textbf{(2.)} 

\vspace{3mm}

\noindent\mybox{\colorbox{lightgray}{\textbf{Statement}}}\hspace{3mm}\hrulefill 

\noindent Every open set is an $F_{\sigma}$ set. 

\vspace{5mm}

\noindent\mybox{\colorbox{lightgray}{\textbf{Proof}}}\hspace{3mm}\hrulefill

Let $\mathcal{O}$ be an open set. Then we can write $$\mathcal{O} = \bigcup\limits_{k=1}^{\infty}(a_k, b_k),$$ where $a_k,b_k$ are extended real numbers satisfying $a_k\leq b_k.$ Note that we can write each $(a_k, b_k)$ as
$$
(a_k, b_k) = \bigcup\limits_{n=1}^{\infty}\left[a_k + \frac{b_k-a_k}{3n}, b_k - \frac{b_k-a_k}{3n}\right].
$$
So we can write 
$$
\mathcal{O} = \bigcup_{k=1}^{\infty}\left(\bigcup\limits_{n=1}^{\infty}\left[a_k + \frac{b_k-a_k}{3n}, b_k - \frac{b_k-a_k}{3n}\right]\right).
$$
A countable union of countable unions is still a countable union. Thus, $\mathcal{O}$ is a countable union of closed sets. Therefore, it is an $F_{\sigma}$ set.

\vspace{3mm}

\hfill$\blacksquare$

\pagebreak

\noindent \textbf{(3.)}

\vspace{3mm}

\noindent\mybox{\colorbox{lightgray}{\textbf{Statement}}}\hspace{3mm}\hrulefill

\noindent Let $A\subset \mathbb{R}$. Then for any $\varepsilon>0$, there exists an open set $\mathcal{O}$ such that $A\subset \mathcal{O}$ and $m^{*}(\mathcal{O})\leq m^{*}(A) + \varepsilon$.

\vspace{5mm}

\noindent\mybox{\colorbox{lightgray}{\textbf{Proof}}}\hspace{3mm}\hrulefill

Let $\varepsilon > 0$. First, suppose $m^{*}(A) = \infty$. Note that $\mathbb{R}$ is open and $A\subset\mathbb{R}$. We have
\begin{align*}
    m^{*}(\mathbb{R}) &= \infty \\
    &\leq \infty \\
    &= \infty + \varepsilon \\
    &= m^{*}(A) + \varepsilon,
\end{align*}
as required. 

Next, suppose $m^{*}(A) < \infty.$

\hfill$\blacksquare$

\pagebreak

\noindent \textbf{(4.)}

\vspace{3mm}

\noindent\mybox{\colorbox{lightgray}{\textbf{Statement}}}\hspace{3mm}\hrulefill

\noindent Suppose $A,B$ are subsets of $\mathbb{R}$ and $m^{*}(A) = 0.$ Then $m^{*}\left(A\cup B\right) = m^{*}(B).$

\vspace{5mm}

\noindent\mybox{\colorbox{lightgray}{\textbf{Proof}}}\hspace{3mm}\hrulefill

Clearly $B\subseteq A\cup B$. Thus, by the monotonicity of outer measure, we have $$m^{*}\left(B) \leq m^{*}(A\cup B).$$ We also have that $$m^{*}(A\cup B) \leq m^{*}(A) + m^{*}(B) = 0 + m^{*}(B) = m^{*}(B),$$ by Proposition 3. Since $m^{*}(A\cup B) \leq m^{*}(B)$ and $m^{*}(B) \leq m^{*}(A \cup B)$, it must be the case that $m^{*}(A\cup B) = m^{*}(B).$ 

\hfill$\blacksquare$

\pagebreak

\noindent \textbf{(5.)}

\vspace{3mm}

\noindent\mybox{\colorbox{lightgray}{\textbf{Statement}}}\hspace{3mm}\hrulefill

\noindent Let $A$ be the set of irrational numbers in $[0,1]$. Then $m^{*}(A) = 1.$

\vspace{5mm}

\noindent\mybox{\colorbox{lightgray}{\textbf{Proof}}}\hspace{3mm}\hrulefill

Let $B$ be the set of rational numbers in $[0,1]$. Then we can write $[0,1] = A \cup B$. Now, $\mathbb{Q}$ is countable, and since a subset of a countable set is countable, it follows that $B$ is countable. Therefore, $m^{*}(B) = 0$. Applying the result of Problem 4 (from the previous page), we see that
$$m^{*}([0,1])=m^{*}(A\cup B) = m^{*}(A).$$ But $m^{*}([0,1])=1.$ Thus, $m^{*}(A) = 1.$

\hfill$\blacksquare$
\end{document}
